% Propietario conservado: /Users/ruben/Documents/New project/output/REVISION_ARGUMENTAL_APERTURAS_CIERRES_20260915/VI/delivery/MANUSCRIPT_VI_ES_EN_COMPLETE/source_es/sections/10_sectores_topologicos.tex
\section{Persistence, fusion structures, and spectral realisations}
\label{vi:vi:sec:topologia}

The atlas and mass operators admit different representations. The
classification of these representations preserves the datum characterising
each: a multiplication rule, a real conjugation, an exchange phase,
or a spectral measure. This section places the topological
sectors and neutral realisations within the same architecture.
Mass coordinates are obtained once the representation and its
dynamics have been determined.

The distinction between these stages is particularly visible in the
areal--volumetric sector. An incidence matrix can fix a fusion ring
and its positive dimension; the energy of an excitation also depends on the
operator acting in the representation. The first determination retains
exact content when several dynamic realisations of
the same ring are studied. The distinction preserves the structural information
preceding the magnitude.

\subsection{The based ring of areal--volumetric incidence}
\label{vi:vi:subsec:fusion-av}

The primitive incidence of the sheets, obtained from the TPK calendar, is
\[
 F_{\mathrm{av}}=\begin{pmatrix}0&1\\1&1\end{pmatrix}.
\]
In the basis $(1,\tau)$, the first column expresses
$\tau\cdot1=\tau$; the other expresses $\tau\cdot\tau=1+\tau$.
The incidence thus induces the based ring
\[
 \mathcal F_{\mathrm{av}}=\mathbb Z[\tau]/(\tau^2-\tau-1),
 \qquad \mathcal F_{\mathrm{av}}^+
 =\{a+b\tau:a,b\in\mathbb Z_{\geq0}\}.
\]
The positive cone records multiplicities. The operation is internal to the
two types and does not use a mass as a coefficient.

\begin{proposition}[Product and positive dimension]
\label{vi:vi:prop:fusion-av}
We have
\[
 (a+b\tau)(c+d\tau)
 =(ac+bd)+(ad+bc+bd)\tau.
\]
There is a unique ring homomorphism $d:\mathcal F_{\mathrm{av}}
\to\mathbb R$ with $d(1)=1$ and $d(\tau)>0$. Its value at $\tau$
is the golden coordinate recognised in the realisation of $F_{\mathrm{av}}$.
\end{proposition}
\begin{proof}
Expanding the product and substituting $\tau^2=1+\tau$ gives the coefficients.
A unital homomorphism is determined by $x=d(\tau)$, with
$x^2=x+1$. Its roots are $(1+\sqrt5)/2$ and $(1-\sqrt5)/2$;
only the first is positive. Evaluation at that root defines a
homomorphism of the quotient. The equation identifies the dimension of this
representation; the HMT regional coordinate has been constructed earlier.
\end{proof}

With $F_0=0$, $F_1=1$, and $F_{n+1}=F_n+F_{n-1}$, the powers satisfy
\[
 \tau^n=F_{n-1}+F_n\tau\qquad(n\geq1).
\]
The initial case is immediate; multiplication by $\tau$ reapplies
the incidence and produces the induction step. The coefficients thus retain
their origin as counts of sheet composition.

A categorical representation adds morphism spaces and an
associator to these multiplicities. A braided representation adds
compatible exchange operators. The material realisation adds
the Hamiltonian, its domains, and the preparation selecting a
spectral sector. If it has a gap $\Delta E>0$ and effective velocity
$c_{\mathrm{eff}}$, its gap-mass coordinate is
\[
 m_{\mathrm{brecha}}=\Delta E\,c_{\mathrm{eff}}^{-2}.
\]
The positive dimension of the ring and the gap are outputs of different
operators. The fusion rule preserves its content without by
itself fixing a universal mass for all its realisations.

\subsection{Dimension and representation type}

In a pointed representation, all simple objects are
invertible. If $X$ is simple and invertible and $d$ is a positive
dimension compatible with duals, then
\[
 1=d(X\otimes X^*)=d(X)d(X^*)=d(X)^2,
\]
and $d(X)=1$. Dimension therefore provides a check on
changes of representation.

\begin{proposition}[Dimension obstruction in the pointed sector]
\label{vi:vi:prop:apuntada}
A simple object of dimension $\varphi$ or $\sqrt2$ cannot be identified,
preserving dimension, with a simple object of a pointed representation.
In particular, the $81$ invertible simple objects of the untwisted centre
based on $\mathbb Z/9\mathbb Z$ do not provide that identification.
\end{proposition}
\begin{proof}
Invertible simple objects have dimension $1$, different from both
dimensions. In the stated abelian untwisted centre, a simple object is
labelled by a group element and a character of its dual. Both
sets have $9$ elements. The $81$ labels are invertible
under multiplication of the group and of the character.
\end{proof}

This census of $81$ has a domain different from the atlas of $81$ cells.
An extension or functor to another codomain specifies the actions
and maps it transports. Areal--volumetric incidence preserves
its exact ring, and the topological realisation preserves the conditions
of its own type.

\subsection{Neutral self-conjugacy and real generators of CAR}
\label{vi:vi:subsec:majorana-categorias}

A real structure on a complex space is an antiunitary operator
$\mathcal C$ with $\mathcal C^2=I$. Its fixed part is a real space.
Commutation of the mass operator with $\mathcal C$ preserves that
real part on its domain. Restricting an evolution to the fixed
part is expressed through commutation of its group with $\mathcal C$;
in a Schrodinger evolution, this condition must be distinguished
from commutation of the Hamiltonian. Both conditions belong to the
material representation of Section~\ref{vi:vi:subsec:antisecciones}.

\begin{proposition}[Antiunitary conjugation and time evolution]
\label{vi:vi:prop:conjugacion-evolucion}
Let $H$ be self-adjoint, $\hbar>0$, and $\mathcal C$ an antiunitary
conjugation preserving $\mathcal D(H)$. For
$U(t)=\exp(-itH/\hbar)$, if
$\mathcal C H\mathcal C^{-1}=\epsilon H$, with
$\epsilon\in\{1,-1\}$, then
\begin{equation}
 \mathcal C U(t)\mathcal C^{-1}=U(-\epsilon t).
 \label{vi:vi:eq:conjugacion-grupo-temporal}
\end{equation}
The group $U(t)$ preserves the fixed part of $\mathcal C$ for every $t$
if and only if it commutes with $\mathcal C$, equivalently if
$\mathcal C H\mathcal C^{-1}=-H$.
If instead $\mathcal C H=H\mathcal C$, a vector fixed by
$\mathcal C$ remains fixed by it throughout the evolution
exactly when it belongs to $\ker H$.
\end{proposition}
\begin{proof}
Antiunitary conjugation sends $i$ to $-i$. Spectral functional
calculus gives
\[
 \mathcal C\exp(-itH/\hbar)\mathcal C^{-1}
 =\exp\bigl(it\mathcal C H\mathcal C^{-1}/\hbar\bigr),
\]
which proves \eqref{vi:vi:eq:conjugacion-grupo-temporal}. If $U(t)$ commutes
with $\mathcal C$, it transports its fixed vectors to fixed vectors.
Conversely, every complex vector is uniquely written as
$x+iy$ with $x,y$ fixed by $\mathcal C$. Preservation of that real
part implies, by linearity of $U(t)$ and antilinearity of
$\mathcal C$, that the two commute. Uniqueness of the self-adjoint
generator of the unitary group then identifies
$\mathcal C H\mathcal C^{-1}$ with $-H$.

In the commuting case $\epsilon=1$, for $\mathcal C\psi=\psi$
we have $\mathcal C U(t)\psi=U(-t)\psi$. The condition that $U(t)\psi$ be fixed
for every $t$ is equivalent to $U(2t)\psi=\psi$ for every $t$. By
spectral calculus, this is equivalent to the spectral measure of $\psi$
being supported on $\{0\}$, that is, to $\psi\in\ker H$.
\end{proof}

Thus an antiunitary symmetry commuting with $H$ reverses time
direction. A realisation through preserved real states uses
a different condition on the generator. This distinction preserves the scope
of algebraic self-conjugacy and allows the dynamics of a Majorana
field to be specified separately.

\begin{proposition}[Self-conjugacy and charge]
\label{vi:vi:prop:carga-real}
Let $Q$ be a self-adjoint charge operator and $\mathcal C$ a
conjugation transporting its domain and satisfying
$\mathcal C Q\mathcal C^{-1}=-Q$. If $\psi\ne0$,
$Q\psi=q\psi$, and $\mathcal C\psi=\psi$, then $q=0$.
\end{proposition}
\begin{proof}
Self-adjointness implies $q\in\mathbb R$. Applying $\mathcal C$
to the eigenvalue equation gives $-Q\psi=q\psi$. Together with the initial equation
this yields $2q\psi=0$, hence $q=0$.
\end{proof}

Neutrality is necessary in this self-conjugacy problem;
the kernel of $Q$ does not by itself specify an antiunitary operator
or its compatibility with the evolution group. The criterion applies to the
neutrino blocks after constructing their operators and mixing.
The fixed dodecaphase word, the real spinor, and the charge
representation thus retain their respective domains.

The fermionic completion provides another realisation. On the
exterior space of the stated modes, $a_j^*$ is exterior
multiplication by the basis vector and $a_j$ is its adjoint. Exterior
antisymmetry gives
\[
 \{a_j,a_k\}=0,\qquad \{a_j,a_k^*\}=\delta_{jk}I.
\]

\begin{proposition}[Self-adjoint generators of CAR]
\label{vi:vi:prop:car-real}
The operators
\[
 \gamma_{2j-1}=a_j+a_j^*,\qquad
 \gamma_{2j}=-i(a_j-a_j^*)
\]
are self-adjoint and satisfy $\{\gamma_r,\gamma_s\}=2\delta_{rs}I$.
\end{proposition}
\begin{proof}
Adjunction fixes each operator. For the same mode,
$a_j^2=(a_j^*)^2=0$ and $a_ja_j^*+a_j^*a_j=I$, so each
square equals $I$. The anticommutator between the two combinations
of the same mode vanishes by expansion. For distinct modes,
all elementary anticommutators vanish.
\end{proof}

These are real generators of a Clifford structure. A topological zero
mode additionally requires the dynamics to realise a localised
zero-energy excitation, separated from the continuum by a gap and
stable under the perturbations considered. The algebraic and spectral
properties are linked through the dynamic realisation.

\subsection{Exchange phases and winding of order six}

For \(n\ge2\), in the order-six reader the exponent sum defines
$w_6:B_n\to\mathbb Z/6\mathbb Z$. The Artin relations
preserve this sum: commutation of distant generators retains
two terms, and the relation involving three generators retains three.

\begin{proposition}[Characters of the winding reader]
\label{vi:vi:prop:devanado-seis}
The unitary characters of the quotient are
\[
 \chi_k(b)=\exp\!\left(\frac{2\pi i k}{6}w_6(b)\right),
 \qquad k=0,1,\ldots,5.
\]
A scalar representation of this form descends to the permutation
group exactly for $k=0$ or $k=3$.
\end{proposition}
\begin{proof}
A cyclic character is determined by the image $q$ of $1$,
with $q^6=1$. To descend to permutations, it must satisfy the
additional relation $\sigma_j^2=1$, which requires $q^2=1$.
Among the sixth roots, only $q=1,-1$ satisfy this condition.
Both satisfy all relations of the permutation group.
\end{proof}

The other four characters preserve winding information.
A material species requires the reader to be realised on a
multisection and the dynamics to preserve its structure. Route
reversal changes $w_6$ to $-w_6$ and transports each character to its
complex conjugate.

\subsection{Stability and visibility of a section}

The positivity proved in
Section~\ref{vi:vi:sec:operador-masa} has a spectral consequence.
For a positive self-adjoint operator $M$, functional calculus gives
$\operatorname{Spec}(M^2)\subset[0,\infty)$. In a realisation
identifying $M^2$ with the squared-mass observable on a
physical space with positive inner product, a persistent eigensection
cannot have a negative eigenvalue of that observable. An unstable
direction of the linearisation of a background belongs to another spectral
problem, studied with its domain and dynamics.

Visibility depends on coupling. If $P_X$ is the support of a
section and $P_a$ is the projector of a channel, $P_aP_X=0$ expresses
orthogonality of supports. A nonzero product expresses possible
incidence; its intensity is determined by the coupling operator.
Two projectors can commute and have nonzero product, so
the commutator does not replace that incidence.

In a neutral block, a section sterile with respect to the weak channel
belongs to the kernel of its coupling. If that kernel is invariant
under the mass operator, the restriction has its own spectral
measure. Atlas depth, neutrality and absence of weak
coupling are distinct properties verified on the same
section. The representational possibilities are preserved without
turning a family label into an external mass.

The organisation thus preserves the order of the argument: incidence
determines multiplicities; lifts determine
conjugations and phases; the representation and material operator
determine the spectral output. Comparison is subsequently applied to
that output, with its units and coupling stated.
